\section{Construcción y verificación del programa}

\subsection{Herramientas y bibliotecas empleadas}

El modelo fue implementado en AnyLogic 8.9.10 (Personal Learning Edition)
combinando dos bibliotecas:

\begin{itemize}
  \item La \textbf{Pedestrian Library}, que representa el desplazamiento
  físico de los pacientes sobre el plano de la guardia, incluyendo la
  ocupación del espacio en la sala de espera.
  \item La \textbf{Process Modeling Library}, que aporta las colas, los
  conjuntos de recursos y las demoras necesarias para representar la
  disciplina de atención.
\end{itemize}

La combinación de ambas responde a una necesidad concreta: la Pedestrian
Library no dispone de colas con prioridad, mientras que la Process Modeling
Library no desplaza peatones por el espacio.

\subsection{Estructura de bloques}

El modelo se organiza en cuatro cadenas:

\begin{table}[H]
\centering
\small
\begin{tabular}{ll}
\toprule
\textbf{Cadena} & \textbf{Secuencia de bloques} \\
\midrule
Ingreso peatonal &
  \texttt{pedSourcePeatonal} $\rightarrow$ \texttt{pedSelectRojoPuerta}
  $\rightarrow$ \texttt{pedServiceSecretaria} \\
Triage &
  \texttt{pedGoToSalaEspera} $\rightarrow$ \texttt{queueTriage} $\rightarrow$
  \texttt{seizeTriage} $\rightarrow$ \texttt{pedGoToTriage} \\
 & $\rightarrow$ \texttt{delayTriage} $\rightarrow$ \texttt{releaseTriage} \\
Consultorios &
  \texttt{queuePrioridad} $\rightarrow$ \texttt{seizeConsultorio} $\rightarrow$
  \texttt{pedGoToConsultorio} \\
 & $\rightarrow$ \texttt{delayConsulta} $\rightarrow$ \texttt{releaseConsultorio} \\
Shock Room &
  \texttt{pedGoToShockRoom} $\rightarrow$ \texttt{queueShock} $\rightarrow$
  \texttt{seizeShock} \\
 & $\rightarrow$ \texttt{delayShock} $\rightarrow$ \texttt{releaseShock} \\
\bottomrule
\end{tabular}
\caption{Organización de la lógica del modelo en AnyLogic.}
\end{table}

La enfermera de triage se modela como un conjunto de recursos
(\texttt{ResourcePool}) y no como un servicio con área propia. Esta decisión
es deliberada: permite que la espera se produzca en la sala común y no dentro
del consultorio de triage, lo que resulta indispensable para medir
correctamente el hacinamiento.

\subsection{Escala temporal}
\label{sec:escala}

La licencia Personal Learning Edition de AnyLogic limita las simulaciones que
emplean la Pedestrian Library a cinco horas de tiempo simulado. Dado que el
turno que se desea representar abarca ocho horas, se introdujo un factor de
escala temporal

\[
k = \frac{480 \text{ min reales}}{300 \text{ min simulados}} = 1{,}6
\]

implementado como un parámetro del modelo. Todos los tiempos de servicio se
dividen por $k$, todas las tasas de arribo se multiplican por $k$, y la
velocidad de desplazamiento de los peatones se multiplica por $k$. Los
diámetros de los peatones y las dimensiones del plano no se modifican, por
tratarse de magnitudes espaciales.

Esta transformación preserva exactamente el comportamiento del sistema de
colas. En efecto, la utilización de un recurso con $c$ servidores

\[
\rho = \frac{\lambda\, E[S]}{c}
\]

permanece invariante, puesto que $\lambda \to k\lambda$ y
$E[S] \to E[S]/k$. Por la ley de Little, $L = \lambda W$, las longitudes de
cola también se conservan. Los tiempos medidos, en cambio, resultan en
minutos simulados y se reconvierten multiplicando por $k$.

\subsection{Verificación del programa}

La verificación se realizó de manera incremental, ejecutando el modelo tras
cada incorporación y resolviendo los errores detectados antes de continuar.
Se documentan a continuación los principales defectos encontrados y
corregidos, por constituir evidencia del proceso de verificación.

\begin{table}[H]
\centering
\small
\begin{tabular}{p{5.4cm}p{8.4cm}}
\toprule
\textbf{Defecto detectado} & \textbf{Corrección aplicada} \\
\midrule
Pacientes detenidos sobre las líneas de destino &
Los bloques \texttt{Delay} no desplazan peatones. Se antepuso un
\texttt{PedGoTo} hacia el área correspondiente. \\
Espera de triage producida dentro del consultorio &
Se reemplazó el servicio con área propia por un \texttt{ResourcePool} con
espera en la sala común. \\
Pacientes sin destino asignado al llegar al bloque de derivación &
La asignación se trasladó del \texttt{On exit} al \texttt{On enter} del
bloque de demora, de modo que el valor quede definido antes de la evaluación
de la condición. \\
Discrepancia entre la salida elegida y la salida efectiva &
Para condiciones que se asignan en el mismo evento en que el peatón se
desplaza se activó la opción \emph{Choice is stochastic or volatile}. \\
Pacientes de ambulancia sin gravedad real asignada &
Se incorporó la asignación explícita en el bloque de origen. \\
Rama de internación sin conexión &
Se corrigió el enlace del puerto correspondiente. \\
\bottomrule
\end{tabular}
\caption{Defectos detectados durante la verificación y sus correcciones.}
\end{table}

Las pruebas de verificación aplicadas fueron:

\begin{enumerate}
  \item \textbf{Conservación de entidades.} La cantidad de pacientes
  ingresados debe igualar la suma de los egresados y los presentes en el
  sistema al cierre.
  \item \textbf{Consistencia de las distribuciones de entrada.} La proporción
  de colores observada debe aproximarse a la teórica de la
  tabla~\ref{tab:colores}.
  \item \textbf{Consistencia de los tiempos de servicio.} La media observada
  de cada servicio debe aproximarse a la media de su triangular.
  \item \textbf{Ausencia de bloqueos.} Ningún paciente debe permanecer
  inmóvil sobre una línea de destino durante la corrida.
  \item \textbf{Coherencia de la disciplina de prioridad.} La espera media de
  los pacientes verdes debe superar a la de los naranjas.
\end{enumerate}

